\begin{exercise}{95}
Let $\mathcal C$ be an ordinary category and $S$ a set of morphisms. Show that
every morphism in $h\mathcal C[S^{-1}]$ can be represented by a zig zag of
morphisms in $\mathcal C$, such that the maps pointing in the wrong direction
are contained in $S$.

\begin{proof}
Use the vertex-preserving localization model of Proposition~2.4.8:
the attachments $\Delta^1\to J$ and the inner-horn fillers in its
construction introduce no new vertices.
Thus the canonical functor $\ell\colon\mathcal C\to h\mathcal C[S^{-1}]$
is bijective on objects.
Let $\mathcal O$ be the word category constructed in
Exercise~\ref{ex:94}.
Sending forward generators to their $\ell$-images and formal inverses to
the corresponding inverses defines a functor
$\mathcal O\to h\mathcal C[S^{-1}]$ under $\mathcal C$.
It is an equivalence by Exercise~\ref{ex:94} and is bijective on objects,
so it is an isomorphism.

Every morphism of $\mathcal O$ is represented by a finite word in forward
morphisms of $\mathcal C$ and formal inverses of elements of $S$.
Read a forward letter in its given direction and an inverse letter
backwards along its corresponding arrow of $S$.
The resulting zig zag has the same endpoints as the given morphism,
and its composite is that morphism under the displayed identification.
The empty word represents an identity.
\end{proof}
\end{exercise}
